%% DESIGN: Computer Modern (default LaTeX) | Academic article-class feel, no color | Serif, publications-heavy
\documentclass[11pt,a4paper]{article}
\usepackage[left=1in,right=1in,top=0.8in,bottom=0.8in]{geometry}
\usepackage[T1]{fontenc}
\usepackage{lmodern}
\usepackage{titlesec,enumitem,hyperref}

\hypersetup{colorlinks=true,linkcolor=black,urlcolor=black}
\pagestyle{empty}

\titleformat{\section}{\large\bfseries\scshape}{}{0em}{}[\vspace{-0.5em}\hrule\vspace{0.3em}]
\titleformat{\subsection}[runin]{\bfseries}{}{0em}{}[.]
\titlespacing{\section}{0pt}{1.2em}{0.5em}

\begin{document}

\begin{center}
{\LARGE\bfseries Dr.\ Morgan Chen}\\[0.3em]
{\large AI Research Scientist}\\[0.3em]
\small morgan.chen@university.edu $\cdot$ (555) 345-6789 $\cdot$ Cambridge, MA\\
\small \href{https://scholar.google.com/citations?user=mchen}{Google Scholar} $\cdot$ \href{https://github.com/morganchen}{GitHub} $\cdot$ \href{https://morganchen.ai}{morganchen.ai}
\end{center}

\section{Research Interests}
Large language models, reasoning and planning in AI systems, multi-agent systems, alignment and safety, efficient inference methods, neurosymbolic approaches.

\section{Education}
\noindent\textbf{Ph.D.\ Computer Science} \hfill Massachusetts Institute of Technology, 2020\\
\indent Thesis: \textit{``Compositional Reasoning in Neural Architectures''}\\
\indent Advisor: Prof.\ Sarah Williams\\[0.4em]
\textbf{M.S.\ Computer Science} \hfill Stanford University, 2016\\[0.4em]
\textbf{B.S.\ Mathematics} (\textit{summa cum laude}) \hfill Carnegie Mellon University, 2014

\section{Selected Publications}
\begin{enumerate}[leftmargin=2em,nosep,label={[\arabic*]}]
\item \textbf{Chen, M.}, Park, J., Williams, S. ``Scaling Laws for Compositional Generalization in Transformers.'' \textit{Proceedings of NeurIPS}, 2024.
\item \textbf{Chen, M.}, Liu, R. ``Efficient Reasoning via Sparse Mixture-of-Experts.'' \textit{Proceedings of ICML}, 2024.
\item Park, J., \textbf{Chen, M.} ``Multi-Agent Debate Improves Factual Accuracy in LLMs.'' \textit{Proceedings of ICLR}, 2023.
\item \textbf{Chen, M.}, Williams, S. ``Neurosymbolic Program Synthesis for Mathematical Reasoning.'' \textit{Proceedings of NeurIPS}, 2022.
\item \textbf{Chen, M.}, et al. ``Compositional Attention Networks for Few-Shot Learning.'' \textit{Proceedings of AAAI}, 2021.
\item \textbf{Chen, M.}, Williams, S. ``Hierarchical Reasoning with Structured Transformers.'' \textit{Journal of Machine Learning Research}, 22(1):1--45, 2021.
\end{enumerate}

\section{Experience}
\noindent\textbf{Senior Research Scientist} \hfill DeepMind, London (Remote) \hfill 2022--Present
\begin{itemize}[leftmargin=2em,nosep]
\item Lead research on reasoning capabilities in large language models (team of 6)
\item Developed novel chain-of-thought distillation method adopted across 3 product teams
\item Published 4 first-author papers at top venues; filed 2 patents
\end{itemize}
\vspace{0.4em}

\noindent\textbf{Research Scientist} \hfill Google Brain, Mountain View, CA \hfill 2020--2022
\begin{itemize}[leftmargin=2em,nosep]
\item Contributed to foundational work on instruction-tuned language models
\item Designed evaluation benchmarks for compositional reasoning (adopted by 50+ papers)
\end{itemize}

\section{Teaching \& Mentoring}
\noindent Guest Lecturer, MIT 6.884 ``Advanced NLP'' \hfill Spring 2023, 2024\\
Mentored 8 PhD interns; 3 papers published from internship work\\
Program Committee: NeurIPS, ICML, ICLR, ACL (2021--present)

\section{Awards \& Honors}
\noindent Outstanding Paper Award, NeurIPS 2024 $\cdot$ MIT EECS Dissertation Award 2020 $\cdot$ NSF Graduate Fellowship 2016 $\cdot$ Phi Beta Kappa

\end{document}
